\chapter{Conclusiones}

Se presentan a continuación las conclusiones del estudio realizado en este trabajo y una propuesta de futuras tareas para su continuación y mejora. 


\section{Conclusión}
 
Este trabajo ha analizado un a serie de pacientes en donde se han usado marcadores del Parkinson para ver su evolución a lo largo del tiempo. También se ha comprobado la validez de los parámetros usados. Con esta información, se ha desarrollado un algoritmo usando el entorno Leaspy en la que se ha estimado la progresión de la enfermedad, primero en un paciente y, posteriormente, en un grupo de 423 individuos. La estimación que se ha realizado ha consistido en la obtención de una curva a largo plazo de un grupo de paciente de entrenamiento con el objetivo de lograr una curva fiable para un paciente nuevo a corto plazo. A lo largo del trabajo se ha valorado la fiabilidad de la estimación calculando las prestaciones de los modelos de progresión (DPMs). \\

El análisis al inicio del estudio se ha conseguido de manera exitosa obteniendo unos resultados favorables y que proporcionan conocimientos sobre la influencia de los biomarcadores en el desarrollo del deterioro cognitivo. Tras esto, la evaluación de la progresión de la enfermedad del Parkinson de un paciente ha sido estimada de manera precisa, asentando la base para el estudio de la cohorte de la PPMI. Sin embargo, tras la implementación del análisis del conjunto de pacientes y el cálculo de las prestaciones de los DPMs se han dado valores que pueden indicar un fallo en la realización del algoritmo puesto que demuestran una baja precisión de la estimación.


\section{Desarrollos futuros}
Tras lo presentado en el apartado anterior, cabe concluir que este estudio no esta finalizado.
Un posible desarrollo para el futuro será el estudio del resto de prestaciones características de los modelos predictivos como el clasificador bayesiano o el área bajo la curva. También, será benificionso calcular la media del NMAE para poder comparar mejor cada modelo e investigar el resultado de los SCORE y encontrar una solución para la mejora de la exactitud del modelo.